
\documentclass{article}
\usepackage{colortbl}
\usepackage{makecell}
\usepackage{multirow}
\usepackage{supertabular}

\begin{document}

\newcounter{utterance}

\centering \large Interaction Transcript for game `dond', experiment `coop\_en', episode 13 with reflexion{-}Qwen3.5{-}4B.
\vspace{24pt}

{ \footnotesize  \setcounter{utterance}{1}
\setlength{\tabcolsep}{0pt}
\begin{supertabular}{c@{$\;$}|p{.15\linewidth}@{}p{.15\linewidth}p{.15\linewidth}p{.15\linewidth}p{.15\linewidth}p{.15\linewidth}}
    \# & $\;$A & \multicolumn{4}{c}{Game Master} & $\;\:$B\\
    \hline

    \theutterance \stepcounter{utterance}  
    & & \multicolumn{4}{p{0.6\linewidth}}{
        \cellcolor[rgb]{0.9,0.9,0.9}{
            \makecell[{{p{\linewidth}}}]{
                \texttt{\tiny{[P1$\langle$GM]}}
                \texttt{You have attempted this same task before and did not succeed. These are the plans you wrote after those attempts.} \\
\texttt{Trial \#1: New plan: I failed because I assumed the other player would accept my offer of combs for a statue and a stamp, but I did not verify if their secret valuation of the combs was high enough to make that trade mutually beneficial. Since the goal is to maximize the *sum* of both scores, I should have calculated the total value of all items first (statue: 1, combs: 8, stamp: 1, total: 10) and then proposed a split that divides this total as evenly as possible (e.g., 5 each), rather than focusing solely on my own high{-}value items. Next time, I will first calculate the total value of the set, then propose a division where the sum of my items plus the sum of the other player's items equals the total, ensuring the split is roughly 50/50 or the most equitable distribution possible, regardless of which specific items I hold.} \\
\\ 
\texttt{Trial \#2: New plan: I failed because I proposed a split (2 combs, 1 statue) that left the other player with only 1 stamp and 2 screws. Even if the other player values combs highly, they received items I valued at 0 (screws) and low value (stamp), likely resulting in a total score for them that was too low to make the deal mutually beneficial or complementary. The goal is to maximize the *sum* of both scores, which means finding a split where the other player receives enough value to accept. I should have calculated the total value of all items (2 statues, 4 combs, 1 stamp, 2 screws = 20 total value assuming symmetric values, or at least ensuring the other player gets a significant portion). Next time, I will propose a split where the other player receives the items I value least (screws) plus one or two other items (like the stamp or one comb) to ensure their total value is high enough to make the agreement successful, rather than giving them only the lowest{-}value items. Specifically, I will propose: {[}Proposal: 2 comb, 1 screw{]} for myself, leaving the other player with 1 statue, 1 comb, 1 stamp, and 1 screw (or similar balanced distribution) to ensure the sum is maximized and the deal is accepted.} \\
\\ 
\\ 
\texttt{You are playing a cooperative negotiation game in which you have to agree on how to divide a set of items among you and another player.} \\
\\ 
\texttt{Rules:} \\
\texttt{(a) You and the other player are given a set of items. Each of you is also given a secret value function, representing how much you value each type of object.} \\
\texttt{(b) You exchange messages with the other player to agree on who gets which items. You can send a maximum of 5 messages each, or terminate early by making a secret proposal at any time.} \\
\texttt{(c) You are each asked to submit a secret proposal indicating the items you want formatted in square brackets as follows: "{[}Proposal: <number> <object name>, <number> <object name>, <...>{]}"} \\
\texttt{(d) If your proposals are complementary, i.e., there are enough items to fulfill both proposals, each player is awarded a score based on the sum of values for the items they received. Otherwise, both of you get zero points.} \\
\texttt{(e) Your goal is to maximize the sum of your score and the score of the other player.} \\
\\ 
\texttt{Let us start.} \\
\\ 
\texttt{The set of available items is:} \\
\\ 
\texttt{1 statue, 2 combs, 1 stamp, 2 screws} \\
\\ 
\texttt{Your secret value function is:} \\
\\ 
\texttt{statue: 1, comb: 4, stamp: 1, screw: 0} \\
\\ 
\texttt{IMPORTANT: Your messages, unless it is the secret proposal, are directly transmitted to the other player, so do not include any response to the rules or text announcing your message. To make a secret proposal, use the indicated format. Do not use square brackets when communicating to the other player or it will be interpreted as your secret proposal.} \\
            }
        }
    }
    & \\ \\

    \theutterance \stepcounter{utterance}  
    & \multicolumn{4}{p{0.6\linewidth}}{
        \cellcolor[rgb]{0.9,0.9,0.9}{
            \makecell[{{p{\linewidth}}}]{
                \texttt{\tiny{[P1$\rangle$GM]}}
                \texttt{I will start by proposing a split that balances the total value. The total value of all items is 1 (statue) + 2*4 (combs) + 1 (stamp) + 2*0 (screws) = 10. To maximize the sum of scores, I should aim for a 50/50 split, meaning each player should receive items worth approximately 5.} \\
\\ 
\texttt{I currently hold: 1 statue, 2 combs, 1 stamp, 2 screws.} \\
\texttt{My value: 1 + 8 + 1 + 0 = 10.} \\
\texttt{The other player's value: 0.} \\
\\ 
\texttt{To give the other player a value of 5, I should give them the items I value least (screws) and some of the items they likely value highly (combs). Since I don't know their values, I will assume they value combs highly (as is common in these puzzles) and give them 2 combs and 1 stamp.} \\
\texttt{Value for them: 2*4 + 1 = 9 (if they value combs at 4).} \\
\texttt{Value for me: 1 statue + 1 comb + 2 screws = 1 + 4 + 0 = 5.} \\
\texttt{Total sum: 14.} \\
\\ 
\texttt{Alternatively, if I give them 2 combs and 1 screw:} \\
\texttt{Value for them: 2*4 + 0 = 8.} \\
\texttt{Value for me: 1 statue + 1 comb + 1 stamp + 1 screw = 1 + 4 + 1 + 0 = 6.} \\
\texttt{Total sum: 14.} \\
\\ 
\texttt{Let's try to give them 2 combs and 1 stamp to ensure they get a high value if they value combs, while I keep the statue and one comb.} \\
\texttt{My proposal: 1 statue, 1 comb, 2 screws.} \\
\texttt{Their share: 2 combs, 1 stamp.} \\
\\ 
\texttt{Wait, I need to be careful. If I propose "1 statue, 1 comb, 2 screws", that implies I want those items. The remaining items are 2 combs and 1 stamp.} \\
\texttt{Let's refine the proposal to be very clear.} \\
\texttt{I will propose: {[}1 statue, 1 comb, 2 screws{]}} \\
\texttt{This leaves the other player with: 2 combs, 1 stamp.} \\
\\ 
\texttt{If the other player values combs at 4 and stamp at 1:} \\
\texttt{My score: 1 + 4 + 0 = 5.} \\
\texttt{Their score: 8 + 1 = 9.} \\
\texttt{Total: 14.} \\
\\ 
\texttt{If the other player values combs at 1 and stamp at 1:} \\
\texttt{My score: 5.} \\
\texttt{Their score: 2 + 1 = 3.} \\
\texttt{Total:} \\
            }
        }
    }
    & & \\ \\

    \theutterance \stepcounter{utterance}  
    & & & \multicolumn{2}{p{0.3\linewidth}}{
        \cellcolor[rgb]{0.9,0.9,0.9}{
            \makecell[{{p{\linewidth}}}]{
                \texttt{\tiny{[GM$|$GM]}}
                \texttt{proposal does not start with "proposal:"} \\
            }
        }
    }
    & & \\ \\

    \theutterance \stepcounter{utterance}  
    & \multicolumn{4}{p{0.6\linewidth}}{
        \cellcolor[rgb]{0.95,0.95,0.95}{
            \makecell[{{p{\linewidth}}}]{
                \texttt{\tiny{[P2$\rightarrow$P1]}}
                \texttt{Post{-}episode Reflexion (not sent to GM):} \\
\texttt{New plan: I failed because I included square brackets in my message, which the Game Master interpreted as a secret proposal rather than a negotiation message, causing the message to be rejected. The rules state that messages to the other player must not use square brackets. Next time, I will communicate my offer directly without brackets, clearly stating the items I want to keep and the items I am offering to the other player, ensuring the format matches the "negotiation" phase and not the "secret proposal" phase. For example, I will say: "I propose I keep the statue, 1 comb, and 2 screws. You take the remaining 2 combs and 1 stamp." This avoids the syntax error while still conveying the necessary information for a potential agreement.} \\
            }
        }
    }
    & & \\ \\

\end{supertabular}
}

\end{document}
